\begin{proof}[Proof of \cref{obj:lem:observable-intervention-moments}]
Fix an experiment in the stated class and integers
\[
k\in\{0,\ldots,6\},\qquad k\le t<T.
\]
The parameter conditions in \cref{obj:def:binary-pomdp-class} give
\[
L=\exp(\zeta)\ge1.
\]

\begin{enumerate}
\item \emph{The population moment.}
By \cref{obj:ass:policy-overlap}, the likelihood ratios satisfy
\[
b(a\mid x)=0\ \Longrightarrow\ e(a\mid x)=0,
\qquad
0\le\rho_j\le L.
\]
For each epoch \(j\), define the class of bounded history variables by
\[
\mathcal D_j
:=
\left\{
D:\ D\text{ is a bounded real }
\mathcal F_j^-\text{-measurable random variable}
\right\}.
\]
Regard any function
\[
f:\mathcal S\longrightarrow\mathbb R
\]
as a column vector. For \(D\in\mathcal D_j\),
\cref{obj:ass:pomdp-kernel,obj:ass:sequential-ignorability,obj:ass:policy-overlap}
give the weighted one-step identity
\[
\begin{aligned}
\mathbb E_b[D\rho_j f(S_{j+1})]
&=
\mathbb E_b\!\left[
D\sum_{a:b(a\mid X_j)>0}
b(a\mid X_j)\frac{e(a\mid X_j)}{b(a\mid X_j)}
\sum_{s'\in\mathcal S}\int_{\mathcal Y}
f(s')\,K(dy,s'\mid S_j,a)
\right]\\
&=
\mathbb E_b[D(P_ef)(S_j)].
\end{aligned}
\]
The zero-behavior cells contribute zero target probability, which justifies completing the action sum. Integrating the terminal reward instead, and integrating an unweighted successor state, respectively, gives
\[
\begin{aligned}
\mathbb E_b[D\rho_jY_j]
&=
\mathbb E_b\!\left[
D\sum_a e(a\mid X_j)
\sum_{s'}\int_{\mathcal Y}y\,K(dy,s'\mid S_j,a)
\right]
=\mathbb E_b[Dr_e(S_j)],\\
\mathbb E_b[Df(S_{j+1})]
&=
\mathbb E_b\!\left[
D\sum_a b(a\mid X_j)
\sum_{s'}\int_{\mathcal Y}f(s')\,K(dy,s'\mid S_j,a)
\right]
=\mathbb E_b[D(P_bf)(S_j)].
\end{aligned}
\]
All these integrations use the joint reward and successor-state law.
% lean: CausalSmith.Stat.PomdpLatentOverlapMinimax.integral_history_mul_ratio_nextState
% lean: CausalSmith.Stat.PomdpLatentOverlapMinimax.integral_history_mul_ratio_reward
% lean: CausalSmith.Stat.PomdpLatentOverlapMinimax.integral_history_mul_behavior_nextState

Define the successive backward reward functions by
\[
f_{e,0}:=r_e,\qquad
f_{e,\ell+1}(s):=\sum_{s'\in\mathcal S}P_e(s,s')f_{e,\ell}(s'),
\qquad \ell\in\mathbb N_0.
\]
Starting with the terminal reward identity and removing the preceding \(k\) ratios backward yields
\[
\mathbb E_b[Z_t(k)]
=
\mathbb E_b\!\left[
\left(\prod_{j=t-k}^{t-1}\rho_j\right)r_e(S_t)
\right]
=
\mathbb E_b[f_{e,k}(S_{t-k})].
\]
At each removal, the remaining earlier ratios form a bounded pre-action-history variable. When \(k=0\), the displayed product is empty and equals one.
% lean: CausalSmith.Stat.PomdpLatentOverlapMinimax.integral_phiwScore_eq_futureFunction_from_state

In the zero-based convention, \cref{obj:ass:stationary-start} gives
\[
\mathbb E_b[f(S_0)]
=
\sum_{s\in\mathcal S}d_b(s)f(s),
\qquad d_bP_b=d_b.
\]
The unweighted one-step identity propagates this equality. Indeed, if it holds at epoch \(j\), then
\[
\begin{aligned}
\mathbb E_b[f(S_{j+1})]
&=\mathbb E_b[(P_bf)(S_j)]\\
&=\sum_s d_b(s)\sum_{s'}P_b(s,s')f(s')\\
&=\sum_{s'}(d_bP_b)(s')f(s')
=\sum_{s'}d_b(s')f(s').
\end{aligned}
\]
Thus, by induction,
\[
\mathbb E_b[f(S_j)]=\sum_s d_b(s)f(s),
\qquad 0\le j<T.
\]
% lean: CausalSmith.Stat.PomdpBinaryhiddenRate.integral_binary_curState_eq_stationary

Finally, the backward recursion agrees with matrix powers:
\[
f_{e,0}=P_e^0r_e,\qquad
f_{e,\ell+1}
=P_ef_{e,\ell}
=P_e(P_e^\ell r_e)
=P_e^{\ell+1}r_e.
\]
Consequently,
\[
\begin{aligned}
\mathbb E_b[Z_t(k)]
&=\sum_s d_b(s)(P_e^kr_e)(s)\\
&=\sum_{s'}\left(\sum_s d_b(s)P_e^k(s,s')\right)r_e(s')\\
&=\sum_{s'}(d_bP_e^k)(s')r_e(s')
=d_bP_e^kr_e.
\end{aligned}
\]
This expression is independent of \(t\), so it gives the stated common moment \(m_k\).
% lean: CausalSmith.Stat.PomdpBinaryhiddenRate.binary_markovOperatorIter_eq_mulVec

\item \emph{The second moment.}
Define the window likelihood-ratio product by
\[
W_{t,k}:=\prod_{j=t-k}^{t}\rho_j.
\]
Sequential assignment and overlap imply
\[
\mathbb E_b[\rho_j\mid\mathcal F_j^-]
=
\sum_{a:b(a\mid X_j)>0}
b(a\mid X_j)\frac{e(a\mid X_j)}{b(a\mid X_j)}
=
\sum_a e(a\mid X_j)
=1.
\]
The ratios at earlier epochs are pre-action-history measurable. Removing the last ratio repeatedly therefore gives
\[
\mathbb E_b[W_{t,k}]
=
\mathbb E_b\!\left[\prod_{j=t-k}^{t-1}\rho_j\right]
=\cdots=
\mathbb E_b[1]
=1.
\]
The last expression is the expectation of the empty product; this also covers \(k=0\).
% lean: CausalSmith.Stat.PomdpLatentOverlapMinimax.integral_phiwWindowProduct_eq_one

There are \(k+1\) factors, so
\[
0\le W_{t,k}\le L^{k+1}.
\]
The reward support supplied by \cref{obj:def:binary-pomdp-class} gives, almost surely,
\[
Z_t(k)^2
=Y_t^2W_{t,k}^2
\le W_{t,k}^2
\le L^{k+1}W_{t,k}.
\]
Taking expectations and using the normalization above proves
\[
\mathbb E_b[Z_t(k)^2]
\le L^{k+1}\mathbb E_b[W_{t,k}]
=L^{k+1}.
\]
% lean: CausalSmith.Stat.PomdpBinaryhiddenRate.score_second_moment_le

\item \emph{Reduction of the separated future score.}
Fix an integer \(h\ge k+1\) with \(t+h<T\), and define
\[
u:=t+h,\qquad
g_{\mathrm{gap}}:=h-k-1\in\mathbb N_0,
\qquad
u-k=t+1+g_{\mathrm{gap}}.
\]
Define the future state function on all of \(\mathcal S\) by
\[
f_{\mathrm{future}}:\mathcal S\longrightarrow\mathbb R,
\qquad
f_{\mathrm{future}}(s)
:=
(P_b^{g_{\mathrm{gap}}}P_e^kr_e)(s).
\]
Both \(1\) and \(Z_t(k)\) belong to \(\mathcal D_{t+1}\). For either choice \(D\in\{1,Z_t(k)\}\), integrate the terminal weighted reward at \(u\), then the \(k\) weighted transitions, and finally the \(g_{\mathrm{gap}}\) unweighted transitions. The one-step identities from the first part give
\[
\begin{aligned}
\mathbb E_b[DZ_u(k)]
&=
\mathbb E_b\!\left[
D\left(\prod_{j=u-k}^{u-1}\rho_j\right)r_e(S_u)
\right]\\
&=
\mathbb E_b[D(P_e^kr_e)(S_{u-k})]\\
&=
\mathbb E_b[D f_{\mathrm{future}}(S_{t+1})].
\end{aligned}
\]
The variable \(D\) remains measurable with respect to every intervening pre-action history. For \(k=0\), the ratio product is empty; for \(g_{\mathrm{gap}}=0\), the unweighted iteration is the identity. In particular,
\[
\begin{aligned}
\mathbb E_b[Z_t(k)Z_u(k)]
&=\mathbb E_b[Z_t(k)f_{\mathrm{future}}(S_{t+1})],\\
\mathbb E_b[Z_u(k)]
&=\mathbb E_b[f_{\mathrm{future}}(S_{t+1})].
\end{aligned}
\]
% lean: CausalSmith.Stat.PomdpLatentOverlapMinimax.integral_phiwScore_mul_eq_futureFunction_full
% lean: CausalSmith.Stat.PomdpLatentOverlapMinimax.integral_phiwScore_eq_futureFunction_full

\item \emph{Oscillation of the future function.}
The kernel is a probability law supported on rewards in \([0,1]\), and the target action probabilities sum to one. Hence, for every \(s\in\mathcal S\),
\[
0
\le r_e(s)
=\sum_a e(a\mid x)\sum_{s'}\int_{\mathcal Y}y\,K(dy,s'\mid s,a)
\le\sum_a e(a\mid x)
=1,
\qquad s=(x,h_{\mathrm{hid}})\in\mathcal X\times\mathcal H.
\]
% lean: CausalSmith.Stat.PomdpBinaryhiddenRate.rewardRegression_unit

For a real function on \(\mathcal S\), define its oscillation by
\[
\operatorname{osc}(f):=\max_{s\in\mathcal S}f(s)-\min_{s\in\mathcal S}f(s).
\]
Thus
\[
\operatorname{osc}(r_e)\le1,
\qquad
0<\alpha=\exp(-1/t_0)\le1.
\]
Both transition matrices are stochastic:
\[
P_b(s,s'),P_e(s,s')\ge0,\qquad
\sum_{s'}P_b(s,s')=\sum_{s'}P_e(s,s')=1.
\]
These identities follow by summing the probability kernel over successor states and then averaging over the corresponding action probabilities.

For \(\mu,\mu'\in\Delta(\mathcal S)\), subtracting the midpoint of the range of \(f\) gives
\[
\begin{aligned}
\left|\sum_s(\mu(s)-\mu'(s))f(s)\right|
&=
\left|\sum_s(\mu(s)-\mu'(s))
\left(f(s)-\frac{\max f+\min f}{2}\right)\right|\\
&\le
\frac{\operatorname{osc}(f)}2
\sum_s|\mu(s)-\mu'(s)|\\
&=
\operatorname{osc}(f)\|\mu-\mu'\|_{\mathrm{TV}}.
\end{aligned}
\]
Applying
\cref{obj:ass:target-contraction,obj:ass:behavior-contraction}
to point masses shows that the total variation distance between any two rows of either matrix is at most \(\alpha\). Therefore,
\[
\begin{aligned}
|(P_ef)(s)-(P_ef)(s')|
&\le
\operatorname{osc}(f)
\|P_e(s,\cdot)-P_e(s',\cdot)\|_{\mathrm{TV}}
\le\alpha\operatorname{osc}(f),\\
|(P_bf)(s)-(P_bf)(s')|
&\le
\operatorname{osc}(f)
\|P_b(s,\cdot)-P_b(s',\cdot)\|_{\mathrm{TV}}
\le\alpha\operatorname{osc}(f).
\end{aligned}
\]
Taking the maximum over the two states and iterating, with the identity operator as the zero-step case, gives
\[
\operatorname{osc}(P_e^kr_e)
\le\alpha^k\operatorname{osc}(r_e)
\le\alpha^k
\le1.
\]
Applying the behavior iteration to this bound yields
\[
\operatorname{osc}(f_{\mathrm{future}})
=
\operatorname{osc}(P_b^{g_{\mathrm{gap}}}P_e^kr_e)
\le
\alpha^{g_{\mathrm{gap}}}\operatorname{osc}(P_e^kr_e)
\le\alpha^{g_{\mathrm{gap}}}.
\]
% lean: Causalean.Mathlib.Probability.FiniteMarkovOscillation.oscillationBound_markovOperatorIter

\item \emph{Centering and the covariance bound.}
The scores are bounded, hence square integrable, and the future function is bounded on the finite state space. The preceding product normalization also gives the absolute first-moment bound
\[
|Z_t(k)|=|Y_t|W_{t,k}\le W_{t,k},
\qquad
\mathbb E_b[|Z_t(k)|]\le\mathbb E_b[W_{t,k}]=1.
\]
% lean: CausalSmith.Stat.PomdpBinaryhiddenRate.score_abs_first_moment_le_one

Define the centering constant by
\[
c_{\mathrm{future}}
:=
\mathbb E_b[f_{\mathrm{future}}(S_{t+1})].
\]
For every \(s\in\mathcal S\), the oscillation bound implies
\[
\begin{aligned}
|f_{\mathrm{future}}(s)-c_{\mathrm{future}}|
&=
\left|\mathbb E_b[
f_{\mathrm{future}}(s)-f_{\mathrm{future}}(S_{t+1})]\right|\\
&\le
\mathbb E_b\!\left[
|f_{\mathrm{future}}(s)-f_{\mathrm{future}}(S_{t+1})|
\right]
\le\alpha^{g_{\mathrm{gap}}}.
\end{aligned}
\]
Consequently,
\[
\begin{aligned}
&\left|
\mathbb E_b[Z_t(k)f_{\mathrm{future}}(S_{t+1})]
-\mathbb E_b[Z_t(k)]c_{\mathrm{future}}
\right|\\
&\qquad=
\left|\mathbb E_b\!\left[
Z_t(k)\{f_{\mathrm{future}}(S_{t+1})-c_{\mathrm{future}}\}
\right]\right|\\
&\qquad\le
\alpha^{g_{\mathrm{gap}}}\mathbb E_b[|Z_t(k)|]
\le\alpha^{g_{\mathrm{gap}}}.
\end{aligned}
\]
% lean: CausalSmith.Stat.PomdpLatentOverlapMinimax.abs_integral_mul_sub_mul_integral_le_oscillation

Expanding the centered product and substituting the two future-score identities gives
\[
\begin{aligned}
&\mathbb E_b\!\left[
\{Z_t(k)-\mathbb E_b[Z_t(k)]\}
\{Z_u(k)-\mathbb E_b[Z_u(k)]\}
\right]\\
&\qquad=
\mathbb E_b[Z_t(k)Z_u(k)]
-\mathbb E_b[Z_t(k)]\mathbb E_b[Z_u(k)]\\
&\qquad=
\mathbb E_b[Z_t(k)f_{\mathrm{future}}(S_{t+1})]
-\mathbb E_b[Z_t(k)]c_{\mathrm{future}}.
\end{aligned}
\]
Its absolute value is therefore bounded by
\[
\alpha^{g_{\mathrm{gap}}}=\alpha^{h-k-1},
\]
as required.
% lean: CausalSmith.Stat.PomdpBinaryhiddenRate.scoreCovariance_le_separated
\end{enumerate}
\end{proof}
